% ═══════════════════════════════════════════════════════════════════════════════
% §8 Discussion paragraphs (two), \input into main.tex at \section{Discussion...}.
% The dangerous-flip headline figure (fig:dflip) now lives in §7.3 (main.tex,
% near its first \cref); these paragraphs reference it backward. Paths relative to
% meeting/ICLR2027/. Citations (blau2018perception, cohen2018distribution,
% geifman2017selective) are resolved in references.bib.
% ═══════════════════════════════════════════════════════════════════════════════

% ---------------------------------------------------------------------------
% §8 DISCUSSION — two paragraphs (\paragraph style, matching sec:discussion)
% ---------------------------------------------------------------------------

\paragraph{Enhancement can suppress malignant cues.}
Figure~\ref{fig:dflip} documents a failure that survives even our strongest
enhancement model. For $10$ of the $74$ correctly-called reference melanomas,
feature-DP \visienhance{} v5 drives the melanoma probability below the decision
line, turning a correct malignant call into a benign one; the case triplets in
panel~(b) show why, as asymmetry, irregular borders, and mottled colour are
smoothed away as though they were capture noise. We read this as the inherent
cost of the perception--distortion trade-off \citep{blau2018perception}: an
operator that makes a degraded image look clean cannot, in general, also
preserve every diagnostically load-bearing high-frequency cue, and a generative
restorer would do worse still by hallucinating plausible-but-absent pathology
\citep{cohen2018distribution}. We therefore suggest that the right response to
``enhancement might hurt this image'' is not a sharper enhancer but a different
action: the residual flips in Figure~\ref{fig:dflip} are precisely the empirical
motivation for the \emph{query-for-retake} channel of our four-channel policy
(Theorem~\ref{thm:agent}). When enhancement is expected to harm rather than help,
the system should decline to enhance and ask for a re-capture, keeping the
harmful transformation out of the diagnostic path entirely.

\paragraph{Severe degradation is a retake signal, not an enhancement target.}
The same logic is borne out at the low-quality extreme. On the severe subset we
observe that enhancement \emph{lowers} diagnostic performance by a margin whose
confidence interval excludes zero ($\Delta\mathrm{AUC}=-0.056$, $95\%$ CI
$[-0.085,\,-0.028]$), with a dangerous-flip rate of $0.46$ (E6). Pushing a
heavily degraded image through any deterministic restorer asks it to invent
detail it cannot recover, and in our measurements this systematically erodes
rather than rescues the diagnosis. Our agent design anticipates this: severe
inputs are routed to \emph{query-for-retake} rather than enhance-then-diagnose,
in the spirit of selective prediction, where abstaining on inputs outside the
reliable operating region is preferable to a confident error
\citep{geifman2017selective}. We therefore read these severe-subset losses not
as a failure of the enhancement module but as positive evidence for triage: some
images should be re-acquired, not repaired, and a system that knows which is
which is safer than one that enhances unconditionally.
